\documentclass[10pt,a4paper]{amsart}
\usepackage[margin=19mm]{geometry}
\usepackage{amsmath,amssymb,mathtools}
\usepackage[scr=rsfs,cal=euler]{mathalfa}
\usepackage{xfrac}
\usepackage[unicode,hidelinks,bookmarks=false]{hyperref}
\newcommand{\ie}{i.e.\ }
\newcommand{\cf}{cf.\ }
\newcommand{\resp}{respectively\ }
\newcommand{\Subsection}{\subsection}
\providecommand{\nobreakdash}{\nobreak\hbox{-}}
\setlength{\emergencystretch}{2em}
\multlinegap0pt
\usepackage{amssymb}
\usepackage{enumerate}
\usepackage{xcolor}
\newtheorem{thm}{Theorem}
\title{On perturbation of Hilbert-Schmidt frames}
\author{Jyoti, Lalit Kumar Vashisht}
\date{Source: arXiv:2511.12019v1 (CC BY 4.0)}
\begin{document}
\maketitle

\noindent\emph{Source and licence.} On perturbation of Hilbert-Schmidt frames. Version arXiv:2511.12019v1 (CC BY 4.0), distributed by arXiv under the Creative Commons Attribution 4.0 International licence, http://creativecommons.org/licenses/by/4.0/, as recorded in the arXiv licence metadata for this exact version and shown on the abstract page https://arxiv.org/abs/2511.12019v1. Full licence notice: \texttt{licenses/arxiv-2511.12019-CC-BY-4.0.txt}.\par\smallskip

\noindent\textit{Editorial scope.} Counted unit: the article's thm thm2 (source0.tex lines 187--198). The article's complete original proof of it is delivered with it (source0.tex lines 199--235, one \texttt{proof} environment of 37 lines). No mathematics is written, repaired or re-derived: the statement and the proof are the article's own lines, in the article's own notation, and no retained line is substituted or re-flowed.  No further source-local context is needed: every cross-reference of every retained line resolves inside the two delivered spans.  Nothing was omitted from any retained span, so the package carries no omission record.  No retained line contains a citation, so the wrapper carries no bibliography and no bibliography entry.  No retained line contains a figure environment, an image-inclusion command or drawing code, so no image is delivered and the package declares no asset dependency.  Editorial formatting only, in addition: an AMS \texttt{amsart} preamble that restates the article's own declarations and macros used by the retained lines, namely the environments \texttt{thm} on separate counters, together with the standard packages those lines need.  The retained environments print in this document as Theorem 1 in order of appearance, whereas the article prints the same items with the numbering of its own full document.  The complete native source of the article as distributed in the arXiv e-print for this exact version is bundled unchanged as \texttt{sources/189168/source0.tex}.
\begin{thm}\label{thm2}
Let  $\{G_i\}_{i \in I}$ be an HS-frame for $\mathcal{H}$ with respect to $\mathcal{K}$ with frame bounds $A_G$, $B_G$, and $\{F_i\}_{i \in I}$ be an HS-Bessel sequence for
$\mathcal{H}$ with respect to $\mathcal{K}$ with Bessel bound $B_F$. Fix any subset $\sigma \subseteq I$ with finite cardinality $|\sigma| \leq N < \infty$ and define the weaving  $\{H_i\}_{i \in I}$ as
\begin{align*}
H_i = \begin{cases}
F_i, & i \in \sigma, \\
G_i, & i \in \sigma^c.
\end{cases}
\end{align*}
Assume that  $\|F_i - G_i\| \leq \epsilon$ for all $i \in I$, for $\epsilon>0$ that satisfies $N\epsilon^2 +2\epsilon\sqrt{N} \sqrt{B_G}<A_G$.
Then,  $\{H_i\}_{i \in I}$ is an HS-frame for $\mathcal{H}$ with frame bounds $\big[A_G-N\epsilon^2 -2\epsilon\sqrt{N} \sqrt{B_G}\big]$ and $(B_F + B_G)$.
\end{thm}

\begin{proof}
It is straightforward to obtain the upper frame condition for the weaving sequence $\{H_i\}_{i \in I}$. Precisely, for any $x \in \mathcal{H}$, we have
\begin{align*}
\sum_{i \in I} \|H_i(x)\|_2^2 =\sum_{i \in \sigma} \|F_i(x)\|_2^2 +\sum_{i \in \sigma^c} \|G_i(x)\|_2^2 \leq \sum_{i \in I} \|F_i(x)\|_2^2+\sum_{i \in I} \|G_i(x)\|_2^2 \leq (B_F + B_G)\|x\|^2.
\end{align*}
Next, to prove the lower frame condition, for each $x \in \mathcal{H}$, we compute
\begin{align}\label{eq2.1}
\sum_{i \in I} \|H_i(x)\|_2^2 &=\sum_{i \in I} \|H_i(x)-G_i(x)+G_i(x)\|_2^2\notag\\
& =\sum_{i \in I}[ H_i(x)-G_i(x)+G_i(x), H_i(x)-G_i(x)+G_i(x)]_{\text{tr}}\notag\\
&= \sum_{i \in I} \|H_i(x)-G_i(x)\|_2^2 + 2 \text{Re} \sum_{i \in I} [ H_i(x)-G_i(x), G_i(x) ]_{\text{tr}} + \sum_{i \in I} \|G_i(x)\|_2^2\notag\\
&= \sum_{i \in \sigma} \|F_i(x)-G_i(x)\|_2^2 + 2 \text{Re} \sum_{i \in \sigma} [ F_i(x)-G_i(x), G_i(x) ]_{\text{tr}} + \sum_{i \in I} \|G_i(x)\|_2^2.
\end{align}
We see that
\begin{align}\label{eq2.2}
&2 \text{Re} \sum_{i \in \sigma} [ F_i(x)-G_i(x), G_i(x) ]_{\text{tr}} \notag\\
& \leq 2 \Big|\sum_{i \in \sigma} [ F_i(x)-G_i(x), G_i(x) ]_{\text{tr}} \Big|\notag\\
& \leq 2 \sum_{i \in \sigma}| [ F_i(x)-G_i(x), G_i(x) ]_{\text{tr}}| \notag\\
& \leq 2 \sum_{i \in \sigma}\| F_i(x)-G_i(x)\|_2\| G_i(x)\|_2 \notag\\
&\leq 2 \big(\sum_{i \in \sigma}\| F_i(x)-G_i(x)\|_2^2\big)^{1/2}\big(\sum_{i \in \sigma}\| G_i(x)\|_2^2\big)^{1/2}\notag\\
& \leq 2 \big(\sum_{i \in \sigma}\| F_i-G_i\|^2\|x\|^2\big)^{1/2}\big(\sum_{i \in I}\| G_i(x)\|_2^2\big)^{1/2}\notag\\
&\leq 2 \big(\sum_{i \in \sigma}\epsilon^2\|x\|^2\big)^{1/2}(B_G\|x\|^2)^{1/2}\notag\\
& \leq 2\epsilon\sqrt{N} \sqrt{B_G}\|x\|^2, \, x \in \mathcal{H},
\end{align}
and
\begin{align}
&\sum_{i \in \sigma} \|F_i(x)-G_i(x)\|_2^2 \leq \sum_{i \in \sigma} \|F_i-G_i\|^2 \|x\|^2  \leq \sum_{i \in \sigma} \epsilon^2 \|x\|^2\leq N\epsilon^2 \|x\|^2, \, x \in \mathcal{H}. \label{eq2.3}
\end{align}
Hence, rewriting equation \eqref{eq2.1} and using equations \eqref{eq2.2} and \eqref{eq2.3}, we get
\begin{align*}
\sum_{i \in I} \|H_i(x)\|_2^2 &=\Big|\sum_{i \in \sigma} \|F_i(x)-G_i(x)\|_2^2 + 2 \text{Re} \sum_{i \in \sigma} [ F_i(x)-G_i(x), G_i(x) ]_{\text{tr}} + \sum_{i \in I} \|G_i(x)\|_2^2\Big|\\
&\geq \Big|\sum_{i \in I} \|G_i(x)\|_2^2\Big|-\Big|\sum_{i \in \sigma} \|F_i(x)-G_i(x)\|_2^2 + 2 \text{Re} \sum_{i \in \sigma} [ F_i(x)-G_i(x), G_i(x) ]_{\text{tr}} \Big|\\
& \geq \sum_{i \in I} \|G_i(x)\|_2^2 - \Big|\sum_{i \in \sigma} \|F_i(x)-G_i(x)\|_2^2\Big|-\Big|2 \text{Re} \sum_{i \in \sigma} [ F_i(x)-G_i(x), G_i(x) ]_{\text{tr}} \Big|\\
& \geq A_G\|x\|^2 -N\epsilon^2 \|x\|^2-2\epsilon\sqrt{N} \sqrt{B_G}\|x\|^2\\
&=\big[A_G-N\epsilon^2 -2\epsilon\sqrt{N} \sqrt{B_G}\big]\|x\|^2, \ x \in \mathcal{H}.
\end{align*}
This completes the proof.
\end{proof}

\end{document}
